\paragraph{Post hoc comparison with the finite-$p$ alternative-workflow results.}
The finite-$p$ results were seen before this protocol was written, so this
comparison is post hoc and descriptive (Table~\ref{tab:posthoc}). Each
alternative-workflow finite-$p$ prevented fraction was within 1.4 percentage
points of the alternative-workflow limit (deviations $+1.31$, $+0.05$ and
$\mn0.73$; 1.4 is a bound) and inside its 95\% interval. In the
original workflow, by contrast, the independent profile's finite-$p$ prevention
(67.6\% at the fixed configuration; 66.7\% after the post hoc common-baseline
interpolation) lies well below its limit of 75.1\% (by 7.6 and 8.4 percentage
points computed from unrounded values), while the clustered (55.9\%) and
imperfect (27.0\%) values lie within 2.3 percentage points of theirs. The leading-order rate
$1000p^2\overline K_h/D^{(0)}$ accounts for only 31--60\% of the measured
absolute control rate in the alternative workflow (25--57\% in the original
workflow), so higher-order terms, including environmental routes and secondary
transmission, are large in absolute terms. In the alternative workflow they
reduced control and policy rates nearly in proportion; that is an observation in
this model, not a guaranteed property. The interval comparison cannot
distinguish the full leading-order limit from the simpler two-opportunity
comparator, which also lies inside every interval.

